% capitulo-1.tex
% Capítulo 1 — Introducción
%
% Citations use the citation keys of the project Zotero library; the
% corresponding entries were exported to references.bib. Claims that would
% require a source not currently available in the library are marked with
% % TODO comments and should be resolved before submission.

\chapter{Introducción}
\label{chap:introduction}

\section{Antecedentes}
\label{sec:background}

Los sistemas robóticos y autónomos modernos dependen cada vez más de la visión
por computador para percibir e interpretar su entorno. Tareas como la
detección, la clasificación y el seguimiento de objetos están dominadas por
las redes neuronales profundas y son computacionalmente demandantes; en
muchas aplicaciones deben ejecutarse de manera continua, en tiempo real o
casi tiempo real, sobre hardware ubicado a bordo del propio robot. La
proliferación de dispositivos móviles y de internet de las cosas ha
desplazado la producción de datos hacia el borde de la red, creando presión
por ejecutar cargas de trabajo de inteligencia artificial junto a los
sensores que generan los datos, en lugar de en centros de datos remotos---un
paradigma consolidado bajo el término de inteligencia en el borde
(\emph{edge intelligence}) \cite{zhou2019, deng2020}. El aprendizaje profundo
en el borde (\emph{edge deep learning}) integra estrechamente recursos
computacionales y fuentes de datos, lo que posibilita la toma de decisiones
en tiempo real en sintonía con el entorno \cite{xu2025}. Delegar la
percepción a servidores remotos suele ser impracticable debido a las
restricciones de latencia, ancho de banda y conectividad y, en algunas
aplicaciones, a consideraciones de privacidad de los datos.
% TODO: the privacy motivation is not explicitly stated in the sources
% currently available in the library; find a dedicated source or remove
% the clause.
Como consecuencia, las cargas de trabajo de visión por computador comúnmente
se ejecutan en plataformas de computación embebidas que deben equilibrar la
capacidad computacional con presupuestos estrictos de energía, térmicos y de
costo \cite{zhou2019, xu2025}. El rango de aplicaciones es amplio y
creciente: drones autónomos para el monitoreo de vida silvestre
\cite{lundquist2026}, detección de defectos en el borde en la agricultura y
la industria \cite{muchtar2025, machado2025}, servicios de ciudades
inteligentes \cite{woo2025} y análisis de imágenes médicas \cite{aalishah2025}.

Entre las plataformas embebidas, la familia NVIDIA Jetson es representativa
de los recursos de computación disponibles para robots y sistemas de
inteligencia artificial en el borde. Estos módulos integran núcleos de CPU
Arm con aceleradores de GPU de NVIDIA en un factor de forma eficiente en
energía y se utilizan ampliamente para visión por computador a bordo e
inferencia de aprendizaje profundo (\emph{deep learning}) \cite{scalcon2024}.
En esta memoria de título se utiliza la NVIDIA Jetson Orin Nano como única plataforma
experimental. Estudios recientes indican que el dispositivo es un candidato
creíble para hospedar exactamente la clase de carga de trabajo estudiada
aquí: Scalcon et al.\ evaluaron la familia Orin para monitoreo de video con
inteligencia artificial y encontraron que la configuración de 15~W de la
Orin Nano es adecuada para aplicaciones de vigilancia, sosteniendo en
promedio 13~FPS para YOLOv4 con DeepStream \cite{scalcon2024}, mientras que
Alqahtani et al., al evaluar de manera comparativa modelos de detección de
objetos en dispositivos de borde heterogéneos, concluyeron que la Orin Nano
ofrece el balance global más favorable entre precisión, latencia y energía
en la mayoría de las familias de modelos \cite{alqahtani2026}. Trabajos
complementarios caracterizan cómo la inferencia de redes neuronales en la
Orin Nano depende de la pila de software---\emph{framework}, precisión
numérica y lenguaje \cite{szegedi2026}--- y de los compromisos de energía y
eficiencia a nivel de sistema \cite{tang2026}. De manera notable para este
trabajo, Szegedi et al.\ reportan que, para la detección de objetos, el
preprocesamiento y el postprocesamiento representan una porción sustancial
del tiempo de ejecución total, especialmente a resoluciones de entrada más
altas \cite{szegedi2026}: una aplicación de visión por computador es más que
su motor de inferencia. Comprende una cadena de etapas---adquisición de
imágenes, preprocesamiento, inferencia y postprocesamiento--- y la forma en
que esas etapas se conectan, dentro de un proceso único o a través de
procesos que se comunican, es una decisión arquitectónica con consecuencias
de rendimiento propias. Esa decisión es el objeto de esta memoria.

En el ámbito del software, los sistemas robóticos han adoptado desde hace
tiempo arquitecturas modulares. El Robot Operating System (ROS), introducido
en 2009, no es un sistema operativo en el sentido tradicional de gestión y
planificación de procesos, sino una capa de comunicación estructurada sobre
el sistema operativo anfitrión \cite{quigley2009}. Su modelo basado en
publicación y sus componentes reutilizables aceleraron demostrablemente la
investigación en robótica \cite{macenski2022}. ROS~2, su segunda generación,
fue rediseñado desde sus bases para cumplir requisitos que ROS~1 no
anticipaba, incluidos el soporte para sistemas de tiempo real y los
despliegues completamente distribuidos y multirobot \cite{macenski2022,
portugal2025}. En términos arquitectónicos, ROS~2 reemplazó el máster
central de ROS~1 por una capa de comunicación descentralizada construida
sobre el estándar Data Distribution Service (DDS), que aporta descubrimiento,
políticas configurables de calidad de servicio (QoS) y primitivas de
transporte adecuadas para sistemas embebidos distribuidos de tiempo real
\cite{macenski2022, maruyama2016, gambo2025}. El \emph{middleware} basado en
DDS se ha extendido ampliamente en la robótica \cite{gambo2025}, y ROS~2 se
utiliza actualmente en dominios tan diversos como vehículos autónomos,
salud, agricultura, aeroespacial y logística \cite{al-batati2026}, con sus
propiedades de tiempo real como objeto de investigación activa
\cite{casini2026}. Con el término del soporte a largo plazo de ROS~1 en mayo
de 2025, la migración de la comunidad a ROS~2 ya no es opcional
\cite{portugal2025}: cómo estructurar aplicaciones en ROS~2 es una pregunta
de diseño presente y urgente para esencialmente todo equipo de este
ecosistema.

En ROS~2, las aplicaciones se descomponen en \emph{nodos}: procesos
independientes que intercambian mensajes tipados a través de canales lógicos
llamados \emph{tópicos}. Este modelo basado en nodos es la fuente de la
célebre modularidad de ROS~2---separación de responsabilidades, desarrollo y
pruebas independientes, y reutilización de componentes entre sistemas
\cite{macenski2022}--- pero también sitúa la comunicación entre procesos en
la ruta crítica de la aplicación. Cada mensaje que cruza una frontera entre
nodos puede incurrir en manejo de mensajes, serialización, copias de
memoria, interacción con el planificador y procesamiento de transporte.
Este sobrecosto de comunicación no es hipotético. Maruyama et al.\
observaron, ya en 2016, que ROS~2 debe convertir los datos para DDS y
abstraer DDS de sus usuarios, y que ello incurre en un sobrecosto adicional
\cite{maruyama2016}. Kronauer et al.\ perfilaron sistemas ROS~2 multinodo y
encontraron que las capas de ROS~2 añaden un sobrecosto de latencia de
hasta 50\% con respecto a la comunicación DDS directa para mensajes
pequeños, con una magnitud que depende de la implementación de DDS, del
tamaño de la carga útil y de la configuración \cite{kronauer2021}. Kamiya
et al.\ rastrearon la latencia a través de las capas internas de la
comunicación de ROS~2 y atribuyeron más de la mitad de la latencia total de
la comunicación predeterminada basada en sockets al segmento de transporte
DDS; reemplazar el transporte por sockets por la comunicación por memoria
compartida redujo la latencia total a menos de la mitad \cite{kamiya2025}.
Luo et al.\ mostraron que incluso la comunicación intraproceso, a menudo
asumida como despreciable, exhibe costos que crecen con el tamaño del
mensaje y con la cantidad de suscriptores \cite{luo2025}. En hardware
embebido estos efectos se amplifican: en una Raspberry Pi~4, Kronauer et al.\
midieron la pérdida de aproximadamente el 30\% de los mensajes para cargas
útiles de 500~KB y un aumento extremo y no lineal de la latencia una vez que
las cadenas superaban unos 13--17 nodos con frecuencias de publicación
superiores a 50~Hz \cite{kronauer2021}. El costo en recursos es igualmente
tangible: componer nodos de ROS~2 en menos procesos ahorró 28\% de CPU y
33\% de RAM para un sistema de navegación completo ejecutándose en un
procesador ARM, y Macenski et al.\ reportan que las cadenas de procesamiento
que transportan mensajes mayores que 1~MB---con imágenes y nubes de puntos
como ejemplos canónicos--- son prácticamente imposibles sin dicha
composición \cite{macenski2023}. Los flujos de video se encuentran entre las
cargas útiles más grandes y frecuentes en los sistemas robóticos, dado que
las plataformas autónomas procesan grandes volúmenes de datos de sensores a
tasas elevadas \cite{kamiya2025}.

La relevancia práctica de estos sobrecostos depende de la relación entre
los costos de comunicación y el presupuesto de cómputo de la plataforma. En
una máquina de clase estación de trabajo, el costo de mover un fotograma
entre procesos puede ser pequeño en comparación con el costo de procesarlo.
En un módulo embebido con recursos limitados que opera bajo una envolvente
de potencia de unos pocos vatios, el mismo costo puede consumir una fracción
sustancial del presupuesto de tiempo por fotograma, reduciendo la tasa de
procesamiento alcanzable o violando los plazos de latencia. El desarrollador
de una aplicación de visión por computador embebida enfrenta, por lo tanto,
un dilema genuino: adoptar la arquitectura modular, bien respaldada e
integrada al ecosistema que ROS~2 promueve, o conservar un diseño monolítico
de proceso único que minimiza la comunicación interna al precio de la
mantenibilidad y las herramientas. Resolver ese dilema con evidencia, en
lugar de con intuición, requiere mediciones controladas en la clase de
hardware donde la tensión es más aguda.

\section{Planteamiento del problema}
\label{sec:problem-statement}

Una aplicación de visión por computador puede implementarse como un proceso
único en el que la adquisición de imágenes, el preprocesamiento, la
inferencia y el postprocesamiento se ejecutan secuencialmente dentro de un
mismo programa. En este enfoque \emph{monolítico}, los datos pasan entre las
etapas de procesamiento como llamadas a funciones ordinarias en memoria, sin
serialización y sin transporte entre procesos. Alternativamente, la misma
aplicación puede descomponerse en nodos de ROS~2 que se comunican mediante
el intercambio de mensajes a través de tópicos. Esta arquitectura
\emph{basada en nodos} proporciona mayor modularidad, pero introduce
fronteras de comunicación entre las etapas de procesamiento, y cada frontera
puede agregar un sobrecosto en tiempo y recursos computacionales
\cite{maruyama2016, kronauer2021, kamiya2025, luo2025}. En una plataforma
embebida con CPU, memoria y energía limitadas, estos costos adicionales
pueden afectar la tasa de procesamiento y la latencia alcanzables de la
aplicación \cite{kronauer2021, macenski2023}. La magnitud del efecto depende
de la carga de trabajo, del grado de particionamiento y del hardware, y por
lo tanto requiere una caracterización experimental en lugar de supuestos.

Esta tensión no es específica de la robótica. En la ingeniería de software
en general, el compromiso entre arquitecturas monolíticas y distribuidas
(de estilo microservicios) se ha estudiado extensamente, y la evidencia
aconseja no asumir que la distribución es gratuita: en una única máquina, se
ha demostrado que un monolito supera en rendimiento a su contraparte basada
en microservicios \cite{blinowski2022}, las diferencias de rendimiento entre
ambos estilos son medibles y dependientes de la carga de trabajo
\cite{tapia2020}, y una revisión crítica reciente incluso identifica una
tendencia de migrar de vuelta desde arquitecturas distribuidas a monolíticas,
motivada en parte por el costo, la complejidad y el rendimiento
\cite{felisberto2024}. Un sistema multinodo de ROS~2 es, en este sentido,
una arquitectura distribuida ejecutada en un único anfitrión, y los propios
mantenedores de ROS~2 reconocen lo que está en juego: el mecanismo de
composición fue creado específicamente para combatir los problemas de
rendimiento del modelo de un proceso por nodo, y al menos un robot comercial
con menos de 60~MB de RAM disponible no habría podido entregar en absoluto
su aplicación de ROS~2 como procesos separados \cite{macenski2023}. Estudios
de pilas completas de conducción autónoma descomponen la latencia de extremo
a extremo en componentes de comunicación, cómputo e inactividad e
identifican las latencias como problemas críticos en sistemas ROS~2 a escala
de producción \cite{betz2023a}, y variantes arquitectónicas como los
despliegues de microservicios contenerizados han sido evaluadas precisamente
porque estos costos de comunicación importan \cite{betz2024}.

Lo que la literatura revisada aún no proporciona, sin embargo, es una
respuesta controlada a la escala de una cadena de procesamiento de visión
por computador embebida. Los estudios de medición existentes caracterizan el
middleware de forma aislada, utilizando cargas útiles sintéticas y cadenas
publicador--suscriptor \cite{maruyama2016, kronauer2021, kamiya2025,
luo2025, zhang2024}; los estudios de mitigaciones como la composición de
nodos evalúan alternativas a la línea base multiproceso en lugar de la
decisión de la línea base misma \cite{macenski2023}; y las evaluaciones de
pila completa, como las de Autoware, confunden innumerables factores
arquitectónicos y algorítmicos \cite{betz2023a, betz2024}. Hasta donde llega
el conocimiento del autor, ningún estudio en la literatura revisada mantiene
constantes la carga de trabajo de visión por computador, el modelo, los
datos de entrada y el hardware mientras varía únicamente la arquitectura de
la aplicación---proceso monolítico frente a cadena de nodos de ROS~2--- y el
grado de particionamiento de la cadena, en una plataforma embebida de clase
Jetson y con la biblioteca cliente rclpy que muchos desarrolladores de
robótica utilizan en la práctica.
% TODO: support the claim that Python/rclpy is widely used among ROS 2
% developers with adoption data (e.g., from the Portugal et al. survey's
% full text or the ROS 2 documentation); no source in the library states
% this explicitly.
El problema que aborda esta memoria es precisamente ese vacío: cuantificar el
compromiso de rendimiento que acompaña a la adopción de una arquitectura
basada en nodos para una carga de trabajo representativa de visión por
computador en una plataforma embebida, bajo condiciones controladas y
reproducibles.

La relevancia de este problema es triple. Primero, es oportuno: con ROS~1
llegando al fin de su vida útil, los equipos están migrando ahora sus
sistemas de percepción a ROS~2 y deben tomar decisiones de arquitectura con
consecuencias duraderas \cite{portugal2025}. Segundo, es de importancia
práctica: la clase de hardware objetivo---ejemplificada por la Jetson Orin
Nano--- es una elección predominante para la visión por computador en el
borde \cite{scalcon2024, alqahtani2026}, y las cargas útiles involucradas,
los flujos de imágenes, se encuentran entre las más demandantes que una
cadena de ROS~2 transporta de manera rutinaria \cite{kamiya2025,
macenski2023}. Tercero, es metodológicamente necesario: la propia comunidad
de robótica llama a una evaluación comparativa de rendimiento estandarizada
y reproducible del cómputo robótico \cite{mayoral-vilches2024}, y las
comparaciones controladas por carga de trabajo de alternativas
arquitectónicas son un ingrediente necesario de ese esfuerzo. La
contribución de esta memoria es una caracterización experimental de los
compromisos de rendimiento asociados a implementar una cadena de
procesamiento de visión por computador como un conjunto de nodos de ROS~2 en
lugar de como un proceso monolítico en una plataforma NVIDIA embebida,
proporcionando orientación práctica para los desarrolladores que deben
decidir si la modularidad de ROS~2 conlleva una penalización significativa
de rendimiento para esta clase de carga de trabajo. La memoria no cuestiona
los beneficios de modularidad de ROS~2, ni intenta establecer una
clasificación universal de rendimiento de los middlewares de robótica.

\section{Preguntas de investigación}
\label{sec:research-questions}

La pregunta de investigación principal de esta memoria es:

\begin{quote}
\emph{¿Cuál es el impacto sobre el rendimiento de particionar una cadena de
procesamiento de visión por computador en nodos rclpy de ROS~2, en
comparación con una implementación equivalente de proceso único en Python,
cuando se ejecuta en una NVIDIA Jetson Orin Nano?}
\end{quote}

Además, se investiga una pregunta de investigación secundaria:

\begin{quote}
\emph{¿Cómo afecta el aumento del grado de particionamiento de la cadena de
procesamiento de visión por computador en nodos de ROS~2 a la latencia, la
tasa de procesamiento y la utilización de recursos?}
\end{quote}

El estudio se ocupa específicamente del compromiso de rendimiento asociado
a la introducción de fronteras de comunicación de ROS~2 en una aplicación
de visión por computador por lo demás equivalente. No se ocupa de identificar
la manera más rápida posible de ejecutar visión por computador con ROS~2:
esa pregunta involucra implementaciones de middleware, configuraciones de
QoS, entornos de despliegue y lenguajes de programación que aquí se
mantienen deliberadamente constantes, como se describe a continuación.

\section{Objetivos}
\label{sec:objectives}

\subsection{Objetivo general}

Caracterizar experimentalmente la diferencia de rendimiento entre una cadena
de procesamiento de visión por computador monolítica en Python e
implementaciones equivalentes en rclpy de ROS~2 sobre una NVIDIA Jetson
Orin Nano.

\subsection{Objetivos específicos}

\begin{itemize}
  \item Desarrollar una cadena de procesamiento de visión por computador
        representativa, capaz de ejecutarse en la Jetson Orin Nano.
  \item Implementar la cadena como una aplicación monolítica de proceso
        único en Python.
  \item Implementar etapas de procesamiento equivalentes como nodos rclpy
        de ROS~2.
  \item Desarrollar múltiples configuraciones de ROS~2 con distintos grados
        de particionamiento de la cadena.
  \item Establecer una metodología reproducible de evaluación comparativa
        de rendimiento (\emph{benchmarking}).
  \item Medir la latencia de extremo a extremo y la tasa de procesamiento.
  \item Medir la utilización de CPU, RAM y GPU.
  \item Medir, de manera opcional, el consumo de energía, si se dispone de
        una medición confiable.
  \item Comparar estadísticamente los resultados entre configuraciones.
  \item Determinar cómo cambia el impacto sobre el rendimiento a medida que
        la cadena de ROS~2 se particiona en mayor grado.
\end{itemize}

\section{Alcance y limitaciones}
\label{sec:scope}

Esta memoria compara dos implementaciones de la misma aplicación de visión
por computador: una implementación monolítica de proceso único en Python y
una implementación multiproceso de ROS~2 que utiliza rclpy. Ambas
implementaciones comparten el mismo lenguaje de programación y el mismo
entorno de ejecución, de modo que el lenguaje de programación sea una
variable controlada y no un factor bajo investigación. El algoritmo de
visión por computador subyacente, el modelo, los datos de entrada, el
hardware y los parámetros de procesamiento relevantes se mantienen
equivalentes entre las implementaciones. La variable independiente principal
es la arquitectura de la aplicación; una variable independiente secundaria
es el grado de particionamiento de ROS~2. Las variables dependientes son la
latencia de extremo a extremo, la tasa de procesamiento, la utilización de
CPU, RAM y GPU y, de manera opcional, el consumo de energía.

Se utiliza una única distribución de ROS~2, una única implementación de DDS
y una única configuración de QoS fija, con un despliegue sin contenedores.
Estas decisiones se tratan como condiciones experimentales controladas y no
como variables bajo investigación.

Los siguientes temas quedan explícitamente fuera del alcance de esta memoria:

\begin{itemize}
  \item Comparación de diferentes lenguajes de programación
  \item Evaluación comparativa de rendimiento de C++ / rclcpp
  \item Nodos componibles de ROS~2
  \item Comparación de diferentes implementaciones de DDS
  \item Comparación de diferentes configuraciones de middleware de ROS~2
  \item Comparación sistemática de políticas de QoS
  \item Evaluación comparativa de rendimiento de Zenoh
  \item Evaluación comparativa de rendimiento de Docker y de la
        contenerización
  \item Comparación de diferentes distribuciones de ROS~2
  \item Comparación de diferentes plataformas de hardware Jetson
  \item Comparación con otros middlewares de robótica
  \item Evaluación comparativa de rendimiento integral de ROS~2 en su
        conjunto
  \item Evaluación de múltiples aplicaciones de visión por computador no
        relacionadas entre sí
\end{itemize}

Estas exclusiones son deliberadas. Este trabajo no pretende responder cuál es la
manera más rápida posible de ejecutar visión por computador con ROS~2, sino
qué costo en rendimiento está asociado al particionamiento de una aplicación
de visión por computador equivalente en nodos de ROS~2 sobre una NVIDIA
Jetson Orin Nano. Los temas excluidos se discuten como posibles trabajos
futuros en el Capítulo~6. Las limitaciones adicionales del estudio,
incluidos el uso de una única plataforma de hardware, una única carga de
trabajo de visión por computador, una configuración fija de DDS y QoS, e
implementaciones en Python, son limitaciones de alcance de un estudio
controlado y no deficiencias, y se discuten junto con las conclusiones.

\section{Organización de la memoria}
\label{sec:organization}

El resto de esta memoria se organiza de la siguiente manera.

\textbf{Capítulo~2 --- Antecedentes y trabajos relacionados} introduce la
visión por computador en plataformas embebidas, la arquitectura de la
plataforma NVIDIA Jetson, la arquitectura de ROS~2 y sus mecanismos de
comunicación, la carga de trabajo de visión por computador seleccionada y
las métricas de rendimiento utilizadas a lo largo de la memoria. Luego revisa
trabajos relacionados sobre \emph{benchmarks} de ROS~2, estudios de
middleware para robótica, benchmarks de visión por computador embebida y
estudios del sobrecosto de comunicación de ROS~2, e identifica el vacío de
investigación que esta memoria aborda.

\textbf{Capítulo~3 --- Metodología} describe la plataforma experimental y el
entorno de software, la cadena de procesamiento de visión por computador y
sus implementaciones monolítica y de ROS~2, las configuraciones
experimentales, el procedimiento de evaluación comparativa de rendimiento,
la definición de cada métrica de rendimiento, los métodos de análisis de
datos y las medidas adoptadas para asegurar la reproducibilidad.

\textbf{Capítulo~4 --- Resultados} presenta los resultados experimentales,
incluidas la tasa de procesamiento, la latencia de extremo a extremo, la
utilización de CPU, memoria y GPU y, si se mide, el consumo de energía,
junto con la comparación de las configuraciones de ROS~2 con distintos
grados de particionamiento.

\textbf{Capítulo~5 --- Discusión} interpreta los resultados, evalúa las
hipótesis, compara los hallazgos con investigaciones previas y discute
observaciones inesperadas y sus posibles explicaciones.

\textbf{Capítulo~6 --- Conclusiones y trabajos futuros} responde las
preguntas de investigación, resume los hallazgos principales, discute las
limitaciones del estudio y esboza direcciones para futuras investigaciones.
